% !TeX root = ../drafts/c-keccak-gkr.tex
\section{Keccak by GKR}\label{annex:keccak}

This annex proves that every \texttt{KECCAK} instruction of an execution (\S\ref{sec:tab-keccak}) computes Keccak-$f[1600]$~\cite{fips202} correctly. It assumes sumcheck and multilinear extensions (\S\ref{sec:prelim}), and recalls what it needs of GKR~\cite{GKR08} and of Keccak.

\subsection{Credits}\label{keccak:credits}

The design is Lev Soukhanov's, after Hashcaster~\cite{Hashcaster}, whose first sumcheck rounds run on bits, and Hashcaster~2~\cite{Hashcaster2}, a GKR for Keccak in which the prover sends a product of packed bits as a polynomial that the verifier checks through a projection~\cite{RMFESkip}. The packing is a reverse multiplication friendly embedding~\cite{CCXY18}. This annex is that design with sixteen bits packed in degree below $24$, the output never committed, and the final claims opened by the VM's commitment.

\subsection{Overview}\label{keccak:overview}

\paragraph{What is proved.} Each row of the \texttt{KECCAK} table is one permutation (a call of Keccak-$f[1600]$). The prover commits the inputs of the permutations, and the main protocol ends holding claims on their outputs. This annex proves that the claimed outputs are Keccak-$f[1600]$ of the committed inputs.

\paragraph{How.} The protocol is a GKR over the $24$ rounds of Keccak-$f[1600]$. Every claim it handles is a weighted sum of the bits of all the states after some round. It starts from one random combination of the output claims and walks back one \emph{layer} per round, each layer turning a claim on the states after a round into a claim on the states before it. After $24$ layers the claim is on the committed inputs, where a short sumcheck reduces it to the $64$ bit slices of the commitment at one point, which ring switching opens (Annex~\ref{annex:rs}).

\paragraph{A layer.} A round is a linear map, then chi, a product of degree two, then the addition of a constant. A claim passes through the linear map for free, by transposing its weights, so a layer works only on chi. The permutations are independent and identical, so the layer handles them sixteen at a time: an embedding packs one bit of each of sixteen permutations into one binary polynomial, so that one product of such polynomials computes sixteen products of bits. The prover sends the claim's weighted sum of chi computed this way, as one polynomial $G$, and the verifier checks it against the claim through a projection. The verifier then evaluates $G$ at a random point, where each pack of sixteen permutations is one field element, and a sumcheck over the packs and the $1600$ bit positions reduces the result to one claim on the states chi was applied to. Transposing the linear map takes that claim to the round's input.

\paragraph{Plan.} \S\ref{keccak:perm} to \S\ref{keccak:cost} describe the protocol when only the inputs are committed, and \S\ref{keccak:recursion} how the recursive verifier evaluates its weights. The VM also commits each state after twelve rounds, which halves the layers; \S\ref{keccak:segments} gives that variant.

\subsection{Keccak-\texorpdfstring{$f[1600]$}{f[1600]}}\label{keccak:perm}

We write Keccak's five maps theta, rho, pi, chi and iota in words: the letter $\fc$ always denotes sumcheck challenges. In this annex, \emph{row} and \emph{column} refer to Keccak's grid below, and we say \emph{table row} and \emph{witness column} for the VM's.

\paragraph{The state.} A state is $1600$ bits, arranged as a $5\times5$ grid of \emph{lanes} of $64$ bits each~\cite{fips202}. The lanes are numbered row by row: lane $\ell<25$ is in row $\lfloor\ell/5\rfloor$ and column $\ell\bmod5$. Bit $z<64$ of a lane is in \emph{slice} $z$, and the bit at \emph{position} $j=z+64\ell$ is bit $z$ of lane $\ell$.

\paragraph{Moving along a row.} The lane $k$ columns along from lane $\ell$, in the same row with the column counted modulo $5$, is
\[
\along{\ell}{k}=5\lfloor\ell/5\rfloor+\bigl((\ell+k)\bmod5\bigr),
\]
so that one column along from lane $3$ is lane $4$, and from lane $4$ it is lane $0$. For a position, $\along{(z+64\ell)}{k}=z+64\,\along{\ell}{k}$ is the same bit of that lane. Moving $k$ along and then $5-k$ along returns to the start, so $j\mapsto\along{j}{k}$ is a bijection with inverse $j\mapsto\along{j}{5-k}$.

\paragraph{A round.} Keccak-$f[1600]$ is $24$ rounds, numbered $R=1,\dots,24$. Round $R$ maps a state $a$ to a state $a'$ in two stages.
\begin{enumerate}
\item The linear map $\klin$: theta, then rho, then pi. Theta adds to each bit the parity of the previous column at the same slice and the parity of the next column at the previous slice, the parity of a column at a slice being the sum of its five bits there. Rho rotates each lane $\ell$ by a fixed amount $\koff{\ell}$: bit $z$ moves to bit $z+\koff{\ell}\bmod64$. Pi permutes the lanes (lane $0$ stays). All three only add or move bits, so $\klin$ is $\Ftwo$-linear.
\item Chi, then iota. Chi adds to each bit the product of the negated next bit along its row and the bit after that. Iota adds the $64$-bit round constant $\krc_R\in\cube{64}$, FIPS~202's $\mathrm{RC}[R-1]$, to lane $0$.
\end{enumerate}
With $b$ the state between the two stages:
\begin{equation}\label{keccak:eq:round}
b=\klin(a),\qquad a'[j]=b[j]+\bigl(1+b[\along{j}{1}]\bigr)\,b[\along{j}{2}]+[\,j<64\,]\,\krc_R[j].
\end{equation}

\subsection{Statement and commitment}\label{keccak:statement}

\paragraph{Permutations.} Each table row is one \emph{permutation}: the instruction reads a state from memory, adds the last padding bit to it (bit $63$ of lane $16$), and applies Keccak-$f[1600]$. The state after that addition is the permutation's \emph{input}. The table has $2^{\kbatch}$ rows, $\kbatch$ being the table's log-height, at least $9$ (\S\ref{sec:arith}), and we index the permutations by $t\in\cube{\kbatch}$, read low bit first. Padding rows are permutations too: an honest prover fills them with the instruction applied to the zero state (\S\ref{sec:e2e-pad}), whose input is the padding bit alone, and the reduction proves every row.

\paragraph{The commitment.} The prover commits the input of every permutation, and no other Keccak state. A lane is stored as an element of $\K$, bit $z$ being the coefficient of $x^z$, and the inputs form one function
\[
\qkeccak:\cube{\kbatch}\times\cube{5}\to\K,
\]
one witness column and one region of the stack (\S\ref{sec:stacking}). Its arguments are a permutation $t$ and a \emph{slot}: slot $\ell<25$ holds input lane $\ell$, and the other seven slots hold zero. So $\qkeccak$ has $\kbatch+5$ variables; the VM's has one more slot variable, $\nkeccak=\kbatch+6$ (\S\ref{keccak:segments}). Its \emph{bit slices} are the Boolean functions $Q_0,\dots,Q_{63}$ with $\qkeccak=\sum_{z<64}Q_z\,x^z$ (Annex~\ref{annex:rs}): bit $z$ of input lane $\ell$ of permutation $t$ is $Q_z(t,\ell)$.

\paragraph{The output claims.} The table's output lanes are virtual columns (\S\ref{sec:tab-keccak}): output lane $\ell$ is the function $\kout{\ell}:\cube{\kbatch}\to\K$ giving lane $\ell$ of Keccak-$f[1600]$ of each committed input, which nobody commits to. The table sumcheck (\S\ref{sec:e2e-unrolled}) ends with one claim on each of them, all at one point $(r_{\mathrm{low}},r_{\mathrm{pack}})\in\E^{4}\times\E^{\kbatch-4}$: $r_{\mathrm{low}}$ at the low four bits of $t$, $r_{\mathrm{pack}}$ at the others. For each $\ell<25$ the prover has sent $\koutv{\ell}\in\E$ and claims $\koutv{\ell}=\mle{\kout{\ell}}(r_{\mathrm{low}},r_{\mathrm{pack}})$.

\paragraph{The statement.} For $0\le e\le24$, let $a_{e}(j;t)\in\Ftwo$ be the bit at position $j$ after $e$ rounds \eqref{keccak:eq:round} of permutation $t$, starting from its committed input. So $a_{0}(z+64\ell;t)=Q_z(t,\ell)$ is committed, the other $a_e$ are determined by the commitment but not committed, and $\kout{\ell}(t)=\sum_{z<64}a_{24}(z+64\ell;t)\,x^z$. The reduction proves
\begin{equation}\label{keccak:eq:statement}
\koutv{\ell}=\mle{\kout{\ell}}(r_{\mathrm{low}},r_{\mathrm{pack}})\quad\text{for all }\ell<25.
\end{equation}
Nothing else is needed: the table's input lanes are read from $\qkeccak$ (lane $16$ as slot $16$ plus $x^{63}$), and the bus binds them to memory (\S\ref{sec:tab-keccak}).

\subsection{Packing sixteen permutations}\label{keccak:rmfe}

All permutations apply chi at the same positions, so the protocol groups them sixteen at a time into \emph{packs}: $t=(i,p)$, with $i\in\cube{4}$ the low four bits of $t$ and $p\in\cube{\kbatch-4}$ the pack. The bits of a pack at one position form a vector of $\Ftwo^{16}$, whose coordinates we index by $i$. The following lemma lets one product of polynomials compute the sixteen products of two such vectors.

\begin{lemma}[Embedding]\label{lem:rmfe}
There are $\Ftwo$-linear maps $\kemb:\Ftwo^{16}\to\Ftwo^{\le23}[X]$ and $\kproj:\Ftwo^{\le46}[X]\to\Ftwo^{16}$, on polynomials of degree at most $23$ and $46$, such that for all $u,v\in\Ftwo^{16}$, with $\wedge$ the coordinatewise product,
\[
\kproj\bigl(\kemb(u)\cdot\kemb(v)\bigr)=u\wedge v.
\]
\end{lemma}

The idea is the Chinese remainder theorem. Modulo a small polynomial $g$, store a few bits in a polynomial $h$ of low degree as its value at $0$, its value at $1$, and its top coefficient (its value at infinity). These three readings multiply when polynomials do. If the degree of $h$ is small enough, the product of two stored polynomials needs no reduction modulo $g$, so its readings are the products of the stored bits. Seven such moduli hold fifteen bits; the sixteenth rides on a multiple of their product.

\begin{proof}
\emph{Readings.} For a polynomial $h$ of degree at most $D$, its readings are $h(0)$, $h(1)$ and its coefficient of $X^{D}$. If $h$ and $h'$ have degree at most $D$ and $D'$, then $hh'$ has degree at most $D+D'$, and each reading of $hh'$, taken with $D+D'$, is the product of the corresponding readings of $h$ and $h'$.

\emph{Moduli.} Take the seven irreducible polynomials below. Being distinct irreducibles, they are pairwise coprime, and their product $F$ has degree $23$. Modulus $g$ holds $n_g$ bits, $15$ in all: $u_0$ goes to $X$, $u_1$ to $X+1$, $u_2,u_3$ and $u_4,u_5$ to the two cubics, and $u_6,u_7,u_8$, $u_9,u_{10},u_{11}$ and $u_{12},u_{13},u_{14}$ to the three quintics, in the order listed.
{\setlength{\arraycolsep}{3pt}
\[
\begin{array}{c|ccccccc}
g & X & X+1 & X^3+X+1 & X^3+X^2+1 & X^5+X^2+1 & X^5+X^3+1 & X^5+X^3+X^2+X+1\\\hline
\deg g & 1 & 1 & 3 & 3 & 5 & 5 & 5\\
n_g & 1 & 1 & 2 & 2 & 3 & 3 & 3
\end{array}
\]}
Modulus $g$ stores its bits in the polynomial $h_g(u)$ of degree at most $n_g-1$ whose first $n_g$ readings are those bits, in order (for $n_g=1$, the constant equal to the bit). It is linear in $u$. Since $2(n_g-1)<\deg g$ ($0<1$, $2<3$, $4<5$), the product $h_g(u)\,h_g(v)$ is already reduced modulo $g$, and by the above its first $n_g$ readings, taken with degree $2(n_g-1)$, are the products $u_kv_k$ of the bits $g$ holds.

\emph{The maps.} By the Chinese remainder theorem there is one polynomial $H_u$ of degree below $23$ with $H_u\equiv h_g(u)\pmod g$ for every $g$, and it is linear in $u$. Put
\[
\kemb(u)=H_u+u_{15}\,F,
\]
of degree at most $23$ and still $\equiv h_g(u)\pmod g$, since $g$ divides $F$. For $P$ of degree at most $46$, let $\kproj(P)$ have, for each modulus $g$, the first $n_g$ readings of $P\bmod g$ taken with degree $2(n_g-1)$, in the bits $g$ holds, and as bit $15$ the coefficient of $X^{46}$ in $P$. Both maps are $\Ftwo$-linear.

\emph{Correctness.} Let $P=\kemb(u)\,\kemb(v)$. Modulo each $g$, $P\equiv h_g(u)\,h_g(v)$, which is already reduced, so $\kproj(P)$ holds the products of the first fifteen bits. In
\[
P=H_uH_v+u_{15}\,FH_v+v_{15}\,FH_u+u_{15}v_{15}\,F^2,
\]
the first three terms have degree at most $45$, and $F^2$ has degree $46$ and leading coefficient $1$, so the coefficient of $X^{46}$ is $u_{15}v_{15}$.
\end{proof}

Write $\kemb_i$ for the embedding of the unit vector at $i\in\cube{4}$ and $\kemb_{\allone}=\sum_i\kemb_i$ for that of the all-ones vector. Then $\kproj(\kemb_{\allone}\cdot\kemb(u))=u$: multiplying by $\kemb_{\allone}$ is how a single embedding is read back. Each bit of $\kproj$ is the parity of a fixed set of coefficients, so it extends to an $\E$-linear form on $\E^{\le46}[X]$, and for weights $\kw\in\E^{16}$ we put
\[
\kproj_{\kw}(G)=\sum_{i\in\cube{4}}\kw_i\,\kproj(G)_i\qquad(G\in\E^{\le46}[X]).
\]

\subsection{Claims}\label{keccak:claims}

All the claims of the reduction have one shape, so that a layer takes one in and gives one out.

\begin{definition}[Claim]\label{def:keccak-claim}
A \emph{claim at depth $d$}, for $0\le d\le24$, is a tuple $(c,\kw,r,\kpw)$ of a value $c\in\E$, \emph{in-pack weights} $\kw\in\E^{16}$, a point $r\in\E^{\kbatch-4}$ and \emph{position weights} $\kpw\in\E^{1600}$. It asserts
\begin{equation}\label{keccak:eq:claim}
c=\sum_{p\in\cube{\kbatch-4}}\ \sum_{i\in\cube{4}}\ \sum_{j<1600}\eq(r,p)\,\kw_i\,\kpw_j\;a_{24-d}\bigl(j;(i,p)\bigr),
\end{equation}
and is \emph{true} if this holds.
\end{definition}

A claim at depth $d$ is a weighted sum of the bits of the states $d$ rounds before the outputs: depth $0$ is the outputs, depth $24$ the committed inputs. The weight of a bit is a product of three factors. Over the packs it is an equality table, as a sumcheck over them leaves. Over the sixteen permutations of a pack it is an arbitrary vector $\kw$, because a layer sets it to the values at a random point of the sixteen unit embeddings (step 2), which form no equality table. Over the positions it is an arbitrary vector $\kpw$, because a layer multiplies it by $\klin^{\top}$ (step 6).

\paragraph{The first claim.} The verifier samples $r_{\mathrm{lane}}\in\E^{5}$ and sets, at depth $0$,
\[
\kw_i=\eq(r_{\mathrm{low}},i),\qquad r=r_{\mathrm{pack}},\qquad \kpw_{z+64\ell}=\eq(r_{\mathrm{lane}},\ell)\,x^{z},\qquad c=\sum_{\ell<25}\eq(r_{\mathrm{lane}},\ell)\,\koutv{\ell}.
\]
The weights $x^z$ reassemble a lane's bits into the lane as an element of $\K$, so by Fact~\ref{fact:eq} the right side of \eqref{keccak:eq:claim} is $\sum_\ell\eq(r_{\mathrm{lane}},\ell)\,\mle{\kout{\ell}}(r_{\mathrm{low}},r_{\mathrm{pack}})$. The first claim is therefore a random combination of the $25$ output claims, and is true if they are.

\subsection{A layer}\label{keccak:layer}

Layer $d$, for $0\le d<24$, takes a claim $(c,\kw,r,\kpw)$ at depth $d$ to a claim at depth $d+1$, through round $24-d$. Let $b=\klin(a_{23-d})$ be the states that round applies chi to, so that by \eqref{keccak:eq:round}
\begin{equation}\label{keccak:eq:chi}
a_{24-d}(j;t)=b(j;t)+\bigl(1+b(\along{j}{1};t)\bigr)\,b(\along{j}{2};t)+[\,j<64\,]\,\krc_{24-d}[j].
\end{equation}
A layer is six steps, each leaving a claim closer to the round's input:
\begin{center}
\begin{tabular}{l@{\qquad}l}
step & claim left after it \\\hline
1. product polynomial & the prover's $G$ is the polynomial \eqref{keccak:eq:product}\\
2. evaluation & $G(\kpt)$, a sum over packs and positions \eqref{keccak:eq:at-s}\\
3. sumcheck over packs & a sum over positions \eqref{keccak:eq:after-packs}\\
4. sumcheck over positions & three evaluations $\kev_0,\kev_1,\kev_2$ of the embedded $b$\\
5. batching & one claim on $b$\\
6. linear map & one claim at depth $d+1$
\end{tabular}
\end{center}

\paragraph{1. Product polynomial.} Embed every pack of $b$: let $S_{p,j}=\kemb\bigl((b(j;(i,p)))_{i\in\cube{4}}\bigr)$ be the embedding of the sixteen bits of pack $p$ at position $j$, and
\[
G_{p,j}=\kemb_{\allone}\,S_{p,j}+\bigl(\kemb_{\allone}+S_{p,\along{j}{1}}\bigr)\,S_{p,\along{j}{2}}.
\]
By linearity $\kemb_{\allone}+S_{p,\along{j}{1}}$ embeds the negated bits, so by Lemma~\ref{lem:rmfe} and $\kproj(\kemb_{\allone}\cdot\kemb(u))=u$, $\kproj(G_{p,j})$ is the sixteen bits of \eqref{keccak:eq:chi} at position $j$, without the round constant. The prover sends the $47$ coefficients of
\begin{equation}\label{keccak:eq:product}
G(X)=\sum_{p\in\cube{\kbatch-4}}\ \sum_{j<1600}\eq(r,p)\,\kpw_j\,G_{p,j}(X)\ \in\E^{\le46}[X],
\end{equation}
the in-pack weights being applied afterwards, by $\kproj_{\kw}$: $\kproj_{\kw}(G)$ is the right side of \eqref{keccak:eq:claim} without the round constant. That constant is public and the same in every permutation, so with $\kwsum=\sum_i\kw_i$ and $\sum_p\eq(r,p)=1$, the verifier checks
\begin{equation}\label{keccak:eq:projection}
\kproj_{\kw}(G)+\kwsum\sum_{z<64}\krc_{24-d}[z]\,\kpw_z\ \qeq\ c.
\end{equation}

\paragraph{2. Evaluation.} The verifier samples $\kpt\in\E$, which reduces the claim that $G$ is \eqref{keccak:eq:product} to the claim that they agree at $\kpt$. At $\kpt$ an embedding is one field element, linear in the bits it embeds:
\[
b^{\kpt}_j(p)=S_{p,j}(\kpt)=\sum_{i}\kw'_i\,b(j;(i,p)),\qquad \kw'_i=\kemb_i(\kpt),\qquad \kwsum'=\sum_i\kw'_i=\kemb_{\allone}(\kpt).
\]
The $\kw'$ are the next claim's in-pack weights. For each position $j$, $b^{\kpt}_j$ is a function on $\cube{\kbatch-4}$, and
\begin{equation}\label{keccak:eq:at-s}
G(\kpt)=\sum_{p\in\cube{\kbatch-4}}\eq(r,p)\sum_{j<1600}\kpw_j\Bigl(\kwsum'\,b^{\kpt}_j(p)+\bigl(\kwsum'+b^{\kpt}_{\along{j}{1}}(p)\bigr)\,b^{\kpt}_{\along{j}{2}}(p)\Bigr),
\end{equation}
whose left side the verifier computes from $G$. This is the heart of the protocol: the sixteen permutations of a pack have become one field element, and the right side is chi again, on field elements instead of bits.

\paragraph{3. Packs.} Steps 3 and 4 are one sumcheck on \eqref{keccak:eq:at-s}, of $\kbatch+7$ rounds, with running claims $\rc{k}$ as in \S\ref{sec:prelim}, starting from $\rc{0}=G(\kpt)$. Its first $\kbatch-4$ rounds bind the pack, each $b^{\kpt}_j$ replaced by its multilinear extension. Their summand is $\eq(r,\cdot)$ times a factor of degree two, so a round costs two $\E$-elements (\S\ref{sec:prelim}). With challenges $\fc^{\mathrm{pack}}\in\E^{\kbatch-4}$, they leave
\begin{equation}\label{keccak:eq:after-packs}
\rc{\kbatch-4}=\eq(r,\fc^{\mathrm{pack}})\sum_{j<1600}\kpw_j\Bigl(\kwsum'\,\mle{b^{\kpt}_j}(\fc^{\mathrm{pack}})+\bigl(\kwsum'+\mle{b^{\kpt}_{\along{j}{1}}}(\fc^{\mathrm{pack}})\bigr)\,\mle{b^{\kpt}_{\along{j}{2}}}(\fc^{\mathrm{pack}})\Bigr).
\end{equation}

\paragraph{4. Positions.} The other $11$ rounds bind the position $j$, written in $\cube{11}$ with the slice in the low six bits and the lane in the high five; every table indexed by position is zero on the $448$ points with lane $25$ or more. The summand is $\kpw$ times a polynomial of degree two in the shifted tables, of degree three in each variable, so a round costs three $\E$-elements. With challenges $\fc^{\mathrm{pos}}\in\E^{11}$, the prover sends, for $k=0,1,2$, the multilinear extension in $j$ of $j\mapsto\mle{b^{\kpt}_{\along{j}{k}}}(\fc^{\mathrm{pack}})$ at $\fc^{\mathrm{pos}}$:
\[
\kev_k=\sum_{j<1600}\eq(\fc^{\mathrm{pos}},j)\,\mle{b^{\kpt}_{\along{j}{k}}}(\fc^{\mathrm{pack}}).
\]
The verifier computes $\mle{\kpw}(\fc^{\mathrm{pos}})=\sum_j\eq(\fc^{\mathrm{pos}},j)\,\kpw_j$ itself and checks the sumcheck's final claim:
\begin{equation}\label{keccak:eq:terminal}
\eq(r,\fc^{\mathrm{pack}})\;\mle{\kpw}(\fc^{\mathrm{pos}})\;\bigl(\kwsum'\,\kev_0+(\kwsum'+\kev_1)\,\kev_2\bigr)\ \qeq\ \rc{\kbatch+7}.
\end{equation}

\paragraph{5. Batching.} The three values are claims on the same states at shifted positions. The verifier samples $\kbeta\in\E$ and combines them into $c'=\kev_0+\kbeta\,\kev_1+\kbeta^{2}\,\kev_2$. Moving each shift onto the weights, $k$ columns along being undone by $5-k$,
\begin{equation}\label{keccak:eq:flat}
c'=\sum_{j<1600}\kpw^{\flat}_j\,\mle{b^{\kpt}_j}(\fc^{\mathrm{pack}}),\qquad
\kpw^{\flat}_j=\eq(\fc^{\mathrm{pos}},j)+\kbeta\,\eq(\fc^{\mathrm{pos}},\along{j}{4})+\kbeta^{2}\,\eq(\fc^{\mathrm{pos}},\along{j}{3}).
\end{equation}
By Fact~\ref{fact:eq}, $\mle{b^{\kpt}_j}(\fc^{\mathrm{pack}})=\sum_{p,i}\eq(\fc^{\mathrm{pack}},p)\,\kw'_i\,b(j;(i,p))$, so $(c',\kw',\fc^{\mathrm{pack}},\kpw^{\flat})$ is a claim of the shape \eqref{keccak:eq:claim} on $b$.

\paragraph{6. Linear map.} Since $b(\cdot;t)=\klin\bigl(a_{23-d}(\cdot;t)\bigr)$ for every $t$, with $\klin^{\top}$ the transpose of $\klin$ extended to $\E^{1600}$,
\[
\sum_{j<1600}\kpw^{\flat}_j\,b(j;t)=\sum_{j<1600}\bigl(\klin^{\top}\kpw^{\flat}\bigr)_j\,a_{23-d}(j;t).
\]
The layer outputs the claim $(c',\,\kw',\,\fc^{\mathrm{pack}},\,\klin^{\top}\kpw^{\flat})$ at depth $d+1$, with no further interaction.

\subsection{The opening}\label{keccak:openings}

Twenty-four layers leave a claim $(c,\kw,r,\kpw)$ at depth $24$, on the committed inputs $a_{0}(z+64\ell;t)=Q_z(t,\ell)$. Grouping its weights by slot and slice, it reads
\begin{equation}\label{keccak:eq:in}
c=\sum_{i\in\cube{4}}\kw_i\sum_{o\in\cube{5}}\sum_{z<64}\kpw^{\mathrm{in}}(o,z)\;\mle{Q_z}\bigl(i,r,o\bigr),\qquad
\kpw^{\mathrm{in}}(\ell,\,z)=\kpw_{z+64\ell},
\end{equation}
with $\kpw^{\mathrm{in}}$ zero on slots $25$ and above. Ring switching needs the $64$ bit slices at one point, but this claim weights the low four bits of $t$ by $\kw$ and the slot by $\kpw^{\mathrm{in}}$, neither an equality table. A sumcheck over these nine variables brings both to a point: four rounds over $i$, with challenges $\fc^{\mathrm{low}}\in\E^{4}$, then five over the slot, with challenges $\fc^{\mathrm{in}}\in\E^{5}$, each of degree two and costing two $\E$-elements. The prover then sends $s_0,\dots,s_{63}\in\E$, its values of the bit slices at the point reached, and the verifier checks the final claim against
\[
\mle{\kw}(\fc^{\mathrm{low}})\sum_{z<64}\mle{\kpw^{\mathrm{in}}}(\fc^{\mathrm{in}},z)\,s_z,
\]
$\kpw^{\mathrm{in}}$ being extended in the slot only. What is left is
\begin{equation}\label{keccak:eq:in-claim}
s_z=\mle{Q_z}\bigl(\fc^{\mathrm{low}},r,\fc^{\mathrm{in}}\bigr)\qquad(0\le z<64),
\end{equation}
the $64$ bit slices of $\qkeccak$ at one point, which ring switching takes (Annex~\ref{annex:rs}).

\subsection{Protocol summary}\label{keccak:summary}

\begin{protocol}[Keccak by GKR]\label{keccak:protocol}
\leavevmode
\begin{enumerate}
\item The verifier samples $r_{\mathrm{lane}}$, which fixes the first claim (\S\ref{keccak:claims}).
\item For $d=0,\dots,23$, layer $d$ (\S\ref{keccak:layer}): the prover sends $G$; the verifier checks \eqref{keccak:eq:projection} and samples $\kpt$; the parties run $\kbatch-4$ sumcheck rounds over packs and $11$ over positions; the prover sends $\kev_0,\kev_1,\kev_2$; the verifier checks \eqref{keccak:eq:terminal} and samples $\kbeta$.
\item The opening (\S\ref{keccak:openings}): four and five sumcheck rounds, then the prover sends $s_0,\dots,s_{63}$.
\end{enumerate}
The commitment scheme is left with \eqref{keccak:eq:in-claim}, its one ring-switched claim (\S\ref{sec:e2e-unrolled}).
\end{protocol}

\subsection{Cost}\label{keccak:cost}

\paragraph{Proof size.} A layer sends $47+2(\kbatch-4)+33+3=2\kbatch+75$ elements of $\E$, and the opening $2\cdot9+64=82$: $2650$ in all for $\kbatch=16$.

\paragraph{Prover.} The work is in step 1 and the first rounds of step 3, and stays on bits as long as possible.
\begin{itemize}
\item The witness is bit-sliced: a machine word holds one position of $64$ permutations, four packs, so building it runs Keccak-$f[1600]$ on $64$ permutations at a time.
\item $G$ is computed on bits, before $\kpt$ exists: a pack and a position cost two carry-less products of binary polynomials of degree at most $23$, $\kemb_{\allone}(S_{p,j}+S_{p,\along{j}{2}})$ and $S_{p,\along{j}{1}}S_{p,\along{j}{2}}$, and field arithmetic starts only when the products meet their weights.
\item After $\kpt$, the prover computes in the basis of the powers of $\kpt$, where a pack's value $b^{\kpt}_j$ is its embedding's bits and multiplying by it is a carry-less product. This basis exists unless $\kpt$ lies in a proper subfield of $\E$; those have fewer than $2^{97}$ elements together, and the prover aborts in that case, with probability below $2^{-95}$ a layer.
\item So the first four rounds of step 3 also run on bits~\cite{Hashcaster}: the sums the prover needs at the points of $\{0,1,\infty\}^4$ are sums of carry-less products of bit-level combinations of packs.
\end{itemize}

\paragraph{Verifier.} A layer's real work is on the $1600$ position weights (\S\ref{keccak:recursion}).

\subsection{Evaluating the weights in the recursion}\label{keccak:recursion}

The native verifier computes the $1600$ position weights of every claim and applies $\klin^{\top}$ to them. The recursive verifier (\S\ref{sec:recursion}) cannot afford this, and never applies $\klin^{\top}$ explicitly.

\paragraph{Factored weights.} Every claim after the first has weights $\klin^{\top}\kpw^{\flat}$, with $\kpw^{\flat}$ as in \eqref{keccak:eq:flat}. A shift along a row moves whole lanes and keeps the slice, so $\kpw^{\flat}$ is a lane table times a slice table:
\[
\kpw^{\flat}_{z+64\ell}=\kpw^{\mathrm{ln}}_{\ell}\,\kpw^{\mathrm{sl}}_{z},\qquad
\kpw^{\mathrm{sl}}_z=\eq(\fc^{\mathrm{pos}}_{<6},z),\qquad
\kpw^{\mathrm{ln}}_\ell=\eq(\fc^{\mathrm{pos}}_{\ge6},\ell)+\kbeta\,\eq(\fc^{\mathrm{pos}}_{\ge6},\along{\ell}{4})+\kbeta^{2}\,\eq(\fc^{\mathrm{pos}}_{\ge6},\along{\ell}{3}).
\]
So the verifier keeps these two factors of $\kpw^{\flat}$, $25$ lane weights and a slice table, instead of $\klin^{\top}\kpw^{\flat}$. The first claim's weights factor too, with slice table $x^z$, and are not transposed.

\paragraph{Where the weights are used.} The verifier needs them in three places, each a pairing $\inner{\klin^{\top}\kpw^{\flat}}{\kpair}=\inner{\kpw^{\flat}}{\klin(\kpair)}$ with a $\kpair$ that factors the same way, $\kpair_{z+64\ell}=\kpair^{\mathrm{ln}}_\ell\,\kpair^{\mathrm{sl}}_z$:
\begin{itemize}
\item $\mle{\kpw}(\fc^{\mathrm{pos}})$ in \eqref{keccak:eq:terminal} of the next layer, where $\kpair$ is the equality table of that layer's $\fc^{\mathrm{pos}}$;
\item the round constants in \eqref{keccak:eq:projection}, where $\kpair$ is supported on lane $0$ and on its slices $0,1,3,7,15,31,63$, the only ones a round constant sets;
\item the final check of the opening, where $\kpair^{\mathrm{sl}}_z=s_z$ and $\kpair^{\mathrm{ln}}_\ell=\eq(\fc^{\mathrm{in}},\ell)$.
\end{itemize}

\paragraph{Pushing $\kpair$ through $\klin$.} Pi permutes lanes, rho rotates the slices of each lane, and theta adds column sums from the same slice and from the previous one, so $\klin(\kpair)$ is a sum of lane tables times rotated slice tables. With $\kto{\ell}$ the lane pi moves lane $\ell$ to, $w_\ell=\ell\bmod5$ its column, and $\kpair^{\mathrm{ln}}_{[w]}=\sum_{k<5}\kpair^{\mathrm{ln}}_{w+5k}$ the sum of $\kpair^{\mathrm{ln}}$ over column $w$ (taken modulo $5$),
\begin{equation}\label{keccak:eq:form}
\inner{\kpw^{\flat}}{\klin(\kpair)}=\sum_{\ell<25}\kpw^{\mathrm{ln}}_{\kto{\ell}}\Bigl(\bigl(\kpair^{\mathrm{ln}}_{\ell}+\kpair^{\mathrm{ln}}_{[w_\ell-1]}\bigr)\,\krot_{\koff{\ell}}+\kpair^{\mathrm{ln}}_{[w_\ell+1]}\,\krot_{\koff{\ell}+1}\Bigr),\qquad
\krot_n=\sum_{z<64}\kpair^{\mathrm{sl}}_z\,\kpw^{\mathrm{sl}}_{(z+n)\bmod64}.
\end{equation}
The slices enter only through the rotation sums $\krot_n$, for $n\in\{\koff{\ell},\koff{\ell}+1\}$. In the first use $\kpair^{\mathrm{sl}}$ is an equality table too, and $\krot_n$ is the multilinear extension of the predicate $[\,z'=z+n\bmod64\,]$ at two points of $\E^{6}$, which a carry ripple over the six bits evaluates with a constant number of multiplications a bit. For the round constants, $\kpair^{\mathrm{sl}}$ has at most seven nonzero entries and the sums are computed directly; for the opening, the $s_z$ are folded against the slice point one bit at a time, for all shifts at once.

\subsection{Optimization for recursion: committing the middle state}\label{keccak:segments}

The proof and the verifier's work are proportional to the number of layers, which is the number of rounds between two committed states, and in the recursion that work is itself proved. The VM therefore also commits each permutation's state after twelve rounds, which doubles $\qkeccak$ and halves the layers. Each layer gains one variable, and the first claim and the openings change.

\paragraph{Commitment and statement.} The $24$ rounds form two \emph{segments} of twelve, segment $m<2$ being rounds $12m+1$ to $12m+12$. Slot $25m+\ell$ of $\qkeccak$ holds lane $\ell$ of segment $m$'s input, the $50$ lanes filling $64$ slots, so there are six slot variables, $\nkeccak=\kbatch+6$. Let $a_{m,e}(j;t)$ be the bit at position $j$ after $e$ rounds of segment $m$, starting from that segment's committed input, so that $a_{m,0}(z+64\ell;t)=Q_z(t,25m+\ell)$ and $\kout{\ell}(t)=\sum_{z<64}a_{1,12}(z+64\ell;t)\,x^z$. The statement is \eqref{keccak:eq:statement} for this $\kout{\ell}$, together with the first segment ending on the committed middle state:
\begin{equation}\label{keccak:eq:chain}
a_{0,12}(z+64\ell;\,t)=Q_z\bigl(t,\,25+\ell\bigr)\quad\text{for all }t\in\cube{\kbatch},\ \ell<25,\ z<64.
\end{equation}

\paragraph{Claims and layers.} Both segments ride one claim, the segment being one more variable next to the pack. A claim at depth $d\le12$ has a point $r\in\E^{\kbatch-3}$, whose last coordinate $r_{\kbatch-4}$ is for the segment, and \eqref{keccak:eq:claim} also sums over $m<2$, with $\eq\bigl(r,(p,m)\bigr)$ in place of $\eq(r,p)$ and $a_{m,12-d}$ in place of $a_{24-d}$. In each segment $m$, layer $d$ goes through round $12(m+1)-d$, the $(12-d)$-th of the segment, for both segments at once. Its round constants enter \eqref{keccak:eq:projection} weighted by $\eq(r_{\kbatch-4},m)$. The sumcheck of steps 3 and 4 gains a round for the segment, after the pack's, $\kbatch+8$ rounds in all: its challenges $\fc^{\mathrm{pack}}$ are then in $\E^{\kbatch-3}$, the last being the next claim's segment coordinate, and \eqref{keccak:eq:terminal} ends on $\rc{\kbatch+8}$. Twelve layers take a claim on both segments' outputs to one on both segments' inputs, which are committed.

\paragraph{The first claim.} The verifier samples $r_{\mathrm{seg}}\in\E$, then $r_{\mathrm{lane}}$, and the first claim, at the point $(r_{\mathrm{pack}},r_{\mathrm{seg}})$, has the weights of \S\ref{keccak:claims} on both segments' outputs. Its part on segment $1$, $\eq(r_{\mathrm{seg}},1)\sum_\ell\eq(r_{\mathrm{lane}},\ell)\,\koutv{\ell}$, comes from the output claims, $\eq(r_{\mathrm{seg}},1)$ being $r_{\mathrm{seg}}$. Its part on segment $0$ is, with $\kmid{\ell}(t)=\sum_{z<64}a_{0,12}(z+64\ell;t)\,x^z$ segment $0$'s output lane $\ell$ as a word, $\eq(r_{\mathrm{seg}},0)\sum_\ell\eq(r_{\mathrm{lane}},\ell)\,\mle{\kmid{\ell}}(r_{\mathrm{low}},r_{\mathrm{pack}})$. No earlier claim supplies it, so the prover sends its value $c_{\mathrm{st}}$, and the first claim's value is $c=c_{\mathrm{st}}+r_{\mathrm{seg}}\sum_\ell\eq(r_{\mathrm{lane}},\ell)\,\koutv{\ell}$.

The first opening below pins $c_{\mathrm{st}}$ to the committed slots $25+\ell$, so the layers test the committed middle state's part plus the output claims' part against what the committed inputs imply. If the committed middle state is segment $0$'s output, the first parts agree by \eqref{keccak:eq:chain}, and the layers test the output claims as in the main text. If it differs in any bit, the two words differ, since the weights $x^z$ make the map from bits to words injective, so the first parts differ as multilinear extensions; and at the random point, $r_{\mathrm{seg}}$ and $r_{\mathrm{lane}}$ being drawn after the output claims are bound, no output claims make up the difference.

\paragraph{The two openings.} After the twelve layers, the parties first open $c_{\mathrm{st}}$, a claim on committed lanes as words:
\begin{equation}\label{keccak:eq:out}
c_{\mathrm{st}}\ \qeq\sum_{o\in\cube{6}}\kpw^{\mathrm{out}}(o)\;\mle{\qkeccak}\bigl(r_{\mathrm{low}},r_{\mathrm{pack}},o\bigr),\qquad
\kpw^{\mathrm{out}}(25+\ell)=\eq(r_{\mathrm{seg}},0)\,\eq(r_{\mathrm{lane}},\ell),
\end{equation}
with $\kpw^{\mathrm{out}}$ zero on the other slots. A degree-two sumcheck over the six slot variables, with challenges $\fc^{\mathrm{out}}\in\E^{6}$, ends with the prover sending $c_{\mathrm{out}}$ and the verifier checking the final claim against $\mle{\kpw^{\mathrm{out}}}(\fc^{\mathrm{out}})\,c_{\mathrm{out}}$. What is left is an ordinary evaluation claim, which joins the main protocol's claim pool:
\begin{equation}\label{keccak:eq:out-claim}
\mle{\qkeccak}\bigl(r_{\mathrm{low}},r_{\mathrm{pack}},\fc^{\mathrm{out}}\bigr)=c_{\mathrm{out}}.
\end{equation}
Then the claim at depth $12$, with point $r\in\E^{\kbatch-3}$, is opened as in \S\ref{keccak:openings}, with $r_{<\kbatch-4}$ in place of $r$, $\kpw^{\mathrm{in}}(25m+\ell,z)=\eq(r_{\kbatch-4},m)\,\kpw_{z+64\ell}$, and six slot rounds instead of five. In \S\ref{keccak:recursion}, its lane table becomes $\kpair^{\mathrm{ln}}_\ell=\sum_{m<2}\eq(r_{\kbatch-4},m)\,\eq(\fc^{\mathrm{in}},25m+\ell)$.

\paragraph{Cost.} The prover sends $c_{\mathrm{st}}$, then $2\kbatch+77$ elements a layer, then $13$ for the first opening (six rounds and $c_{\mathrm{out}}$) and $84$ for the second (ten rounds and the $s_z$): about $1400$ in all for $\kbatch=16$, against $2650$, and the verifier runs half as many layers.

\paragraph{Other segment lengths.} The prover crate \texttt{keccak\_gkr} takes any segment length dividing $24$; the VM's, \texttt{leanvm\_core::keccak::DEPTH}, is $12$.
